%% classification.tex — §5–6 of the standalone paper
%% The G/L/C/M classification and shadow tables

%% =====================================================================
%%  §5.  The Four Shadow Classes
%% =====================================================================

\section{The Four Shadow Classes}\label{sec:four-classes}

The single-line dichotomy (Theorem~\ref{thm:single-line-dichotomy})
partitions the space of modular Koszul algebras into four shadow classes,
determined by two invariants of the shadow metric
$Q_L(t) = (2\kappa + 3\alpha t)^2 + 2\Delta\, t^2$:
the cubic coefficient~$\alpha = S_3$ and the critical
discriminant $\Delta = 8\kappa S_4$.

\begin{definition}[Shadow classes]\label{def:shadow-classes}
Let $\cA$ be a modular Koszul algebra with shadow metric $Q_L$ on a
one-dimensional primary line $L$.  The \emph{shadow class} of $\cA$
along $L$ is:
\[
\renewcommand{\arraystretch}{1.4}
\begin{array}{c|cc|cc}
\textbf{Class} & \alpha & \Delta & r_{\max} & Q_L \\
\hline
G\;\text{(Gaussian)} & 0 & 0 & 2 & (2\kappa)^2 \\
L\;\text{(Lie/tree)} & \neq 0 & 0 & 3 & (2\kappa + 3\alpha t)^2 \\
C\;\text{(Contact)} & 0 & 0^* & 4 & (2\kappa)^2\;\text{(primary line)} \\
M\;\text{(Mixed)} & \neq 0 & \neq 0 & \infty & \text{irreducible quadratic}
\end{array}
\]
Here $r_{\max}$ is the \emph{shadow depth}: the largest arity~$r$
for which the shadow coefficient $S_r \neq 0$.
${}^*$Class~C has $\Delta = 0$ on the primary line (where
the single-line dichotomy applies), but the quartic contact
invariant $Q^{\mathrm{contact}} \neq 0$ lives on a charged
stratum; see Definition~\ref{def:class-C}.
\end{definition}

The classification is exhaustive: every entry in Table~\ref{tab:master-shadow}
falls into exactly one class, and every class is realized by at least
one standard family.  The mechanism is algebraic: the shadow generating
function $H(t) = t^2\sqrt{Q_L(t)}$ is algebraic of degree~$2$
(Theorem~\ref{thm:riccati-algebraicity}), and its Taylor coefficients
$S_r$ are computed from the convolution recursion
$H(t)^2 = t^4 Q_L(t)$, where $H(t) = \sum_{r \geq 2} r\, S_r\, t^r$
and $S_r = a_{r-2}/r$ with $a_n$ the Taylor coefficients of $\sqrt{Q_L}$.
When $Q_L$ is a perfect square ($\Delta = 0$), the square root
terminates; when $Q_L$ is irreducible ($\Delta \neq 0$), the Taylor
expansion is infinite.

\medskip

\subsection{Class G: Gaussian ($r_{\max} = 2$)}

\begin{definition}\label{def:class-G}
A modular Koszul algebra $\cA$ is \emph{Gaussian} if
$S_r(\cA) = 0$ for all $r \geq 3$.  Equivalently,
$\alpha = 0$ and $\Delta = 0$, so the shadow metric
$Q_L(t) = 4\kappa^2$ is constant and the shadow tower
terminates at the scalar curvature.
\end{definition}

\noindent\textbf{Examples.}  Heisenberg $\cH_k$ with $\kappa = k$;
lattice VOAs $V_\Lambda$ with $\kappa = \operatorname{rank}(\Lambda)$;
free fermion $\psi$ with $\kappa = 1/4$;
$bc$ ghosts with $\kappa = -(6\lambda^2 - 6\lambda + 1)$.

\noindent\textbf{Characterization.}
Class~G algebras are freely generated as chiral algebras.  Their bar
complex is quadratic (the differential has no higher-order terms),
and the OPE on each primary line is abelian.  The full obstruction
is captured by a single scalar:
\[
\Theta_\cA = \kappa(\cA) \cdot \eta \otimes \Lambda,
\]
where $\eta$ is the cyclic direction and $\Lambda \in H^*(\bar{\cM}_g)$
is the tautological class.  All genus-$g$ data is controlled by
$F_g(\cA) = \kappa(\cA) \cdot \lambda_g^{\mathrm{FP}}$.

\medskip

\subsection{Class L: Lie/tree ($r_{\max} = 3$)}

\begin{definition}\label{def:class-L}
A modular Koszul algebra $\cA$ is \emph{Lie/tree} if
$\alpha \neq 0$ and $\Delta = 0$, so that $S_3 \neq 0$
but $S_r = 0$ for all $r \geq 4$.  The shadow metric is
a perfect square:
$Q_L(t) = (2\kappa + 3\alpha t)^2$.
\end{definition}

\noindent\textbf{Examples.}  Affine Kac--Moody $V_k(\fg)$ for every
simple Lie algebra $\fg$ and every non-critical level
$k \neq -h^\vee$: type~$A$ ($\mathfrak{sl}_N$),
types $B$, $C$, $D$ (orthogonal and symplectic),
exceptional types $E_6$, $E_7$, $E_8$, $G_2$, $F_4$.
For all of these:
\[
\kappa\bigl(V_k(\fg)\bigr)
= \frac{\dim(\fg)\,(k + h^\vee)}{2h^\vee}\,.
\]

\noindent\textbf{Characterization.}
The nonzero cubic shadow $S_3 = \alpha$ arises from the Lie bracket
(structure constants $f_{abc}$), while the Jacobi identity forces the
quartic shadow $S_4 = 0$ on the primary line.  The $r$-matrix
$r(z) = \Omega_\fg/z$ (single simple pole; the bar construction absorbs
one power of the OPE pole) encodes all nonlinear structure.  The tree-level
graph sum terminates at trivalent vertices.

\medskip

\subsection{Class C: Contact ($r_{\max} = 4$)}

\begin{definition}\label{def:class-C}
A modular Koszul algebra $\cA$ is \emph{contact} if
$\Delta = 0$ on every primary line (so the single-line
dichotomy gives $r_{\max}|_L \leq 3$ on each line), but there
exists a charged stratum direction with $S_4 \neq 0$
(the quartic contact invariant), and the tower terminates at
arity~$4$: $S_r = 0$ for $r \geq 5$.
\end{definition}

\noindent\textbf{Mechanism.}
Class~C escapes the single-line dichotomy (which gives $r_{\max} \in
\{2, 3, \infty\}$) via \emph{stratum separation}: the quartic contact
invariant $Q^{\mathrm{contact}}$ lives on a charged stratum whose
self-bracket exits the cyclic deformation complex at arity~$5$, by
rank-one rigidity.  On the weight-changing primary line, the tower is
abelian ($\alpha = 0$, $S_4 = 0$); the quartic contact is visible
only on the full two-dimensional deformation space.

\noindent\textbf{Examples.}
The $\beta\gamma$ system at conformal weight~$\lambda$, with
\[
\kappa(\beta\gamma_\lambda)
= 6\lambda^2 - 6\lambda + 1
= \tfrac{1}{2}\,c(\beta\gamma_\lambda),
\]
and its Koszul dual, the $bc$ system at the same weight.
At the standard weight $\lambda = 0$ or $\lambda = 1$:
$c = 2$ and $\kappa = 1$.

\medskip

\subsection{Class M: Mixed ($r_{\max} = \infty$)}

\begin{definition}\label{def:class-M}
A modular Koszul algebra $\cA$ is \emph{mixed} if
$\alpha \neq 0$ and $\Delta \neq 0$, so the shadow metric
$Q_L$ is an irreducible quadratic in~$t$ and the shadow tower
does not terminate.
\end{definition}

\noindent\textbf{Examples.}
Virasoro $\mathrm{Vir}_c$ with $\kappa = c/2$;
$\cW$-algebras $\cW_N$ with $\kappa = (H_N - 1)\,c$ where
$H_N = \sum_{j=1}^{N} 1/j$;
Drinfeld--Sokolov reductions of affine Kac--Moody algebras.

\noindent\textbf{Characterization.}
The shadow growth rate
\[
\rho(\cA) = \frac{\sqrt{9\alpha^2 + 2\Delta}}{2|\kappa|}
\]
controls the asymptotic behaviour of $S_r$: the coefficients satisfy
$S_r \sim C(\cA)\, \rho(\cA)^r\, r^{-5/2}\, \cos(r\theta + \phi)$
for a branch-point argument~$\theta$ determined by the shadow metric.
The critical central charge $c^*$ where $\rho(\mathrm{Vir}_c) = 1$ is the
unique positive real root of the cubic
\begin{equation}\label{eq:critical-cubic}
5c^3 + 22c^2 - 180c - 872 = 0,
\qquad
c^* \approx 6.125\,.
\end{equation}
For $c > c^*$ the tower converges ($\rho < 1$), and for $c < c^*$ it
diverges ($\rho > 1$).  At the self-dual point $c = 13$:
$\rho(\mathrm{Vir}_{13}) \approx 0.467$.

The Virasoro shadow tower has explicitly computed coefficients through
arity~$10$:
\begin{align}
S_2 &= \tfrac{c}{2}\,, &
S_3 &= 2\,, &
S_4 &= \frac{10}{c(5c+22)}\,,
\label{eq:vir-S234} \\
S_5 &= \frac{-48}{c^2(5c+22)}\,, &
S_6 &= \frac{80(45c+193)}{3c^3(5c+22)^2}\,, &
S_7 &= \frac{-2880(15c+61)}{7c^4(5c+22)^2}\,.
\label{eq:vir-S567}
\end{align}

\medskip

\begin{remark}[$\mathcal{N}=2$ superconformal algebra]
\label{rem:n2-sca}
The $\mathcal{N}=2$ superconformal algebra at level~$k$ has
central charge $c = 3k/(k+2)$ and modular characteristic
$\kappa = (k+4)/4$.  The Kazama--Suzuki coset realization
gives $\kappa = \kappa(\widehat{\mathfrak{sl}}_2 \text{ at }k)
+ \kappa(\text{fermions}) - \kappa(U(1))
= 3(k{+}2)/4 + 1/2 - (k/2{+}1) = (k+4)/4$.
Under Koszul duality, $c + c' = 6$ (from the $\mathfrak{sl}_2$
involution $k \mapsto -k - 4$).  The algebra is class~$\mathbf{M}$
along the stress-tensor primary line, since the Virasoro subalgebra
has $\alpha = 2$ and $\Delta \neq 0$.
\end{remark}

\begin{remark}[Depth and Koszulness]\label{rem:depth-vs-koszulness}
Shadow depth $r_{\max}$ does \emph{not} characterize chiral Koszulness.
All four classes---Gaussian ($r_{\max} = 2$), Lie ($r_{\max} = 3$),
contact ($r_{\max} = 4$), and mixed ($r_{\max} = \infty$)---consist
of chirally Koszul algebras.  Shadow depth classifies \emph{complexity}
within the Koszul world: the operadic complexity conjecture
identifies $r_{\max}(\cA)$ with the $A_\infty$-depth and
the $L_\infty$-formality level.
\end{remark}

\begin{remark}[Galois splitting]\label{rem:galois-splitting}
The four-class partition has a Galois-theoretic interpretation.
The shadow metric $Q_L(t) = 4\kappa^2 + 12\kappa\alpha\,t
+ (9\alpha^2 + 2\Delta)\,t^2$ has discriminant
$\operatorname{disc}(Q_L) = -32\kappa^2\Delta$.
Over the splitting field:
\begin{itemize}
\item Class~$\mathbf{G}$: $Q_L = (2\kappa)^2$ is a square.
  The ``splitting field'' is $\mathbb{Q}$ itself.
  No extension; trivial Galois group.
\item Class~$\mathbf{L}$: $Q_L = (2\kappa + 3\alpha t)^2$ is
  a square with nontrivial linear factor.  The splitting field
  is still~$\mathbb{Q}$.  The single zero $t_0 = -2\kappa/(3\alpha)$
  is rational.
\item Class~$\mathbf{M}$: $\operatorname{disc}(Q_L) < 0$ at
  positive~$c$.  The zeros of $Q_L$ are a conjugate pair
  $t_0, \bar{t}_0$ in $K_L = \mathbb{Q}(\sqrt{\operatorname{disc}})$,
  an imaginary quadratic extension.  The Galois group
  $\operatorname{Gal}(K_L/\mathbb{Q}) = \mathbb{Z}/2$ acts by
  $t_0 \mapsto \bar{t}_0$, and the shadow tower inherits the
  oscillatory interference from the two conjugate branch points.
\end{itemize}
The passage from class~$\mathbf{L}$ to class~$\mathbf{M}$
is a field extension: the shadow field enlarges from $\mathbb{Q}$
to an imaginary quadratic $K_L$.  The critical discriminant
$\Delta$ is the obstruction to the extension being trivial.
\end{remark}

\begin{remark}[Total depth decomposition]\label{rem:total-depth}
The total shadow depth decomposes as
$d(\cA) = 1 + d_{\mathrm{arith}} + d_{\mathrm{alg}}$
(Theorem~\ref{thm:depth-decomposition}), where the algebraic depth
$d_{\mathrm{alg}} \in \{0, 1, 2, \infty\}$ and depths $\geq 5$ are
purely arithmetic (cusp forms contribute beyond the algebraic engine).
\end{remark}


%% =====================================================================
%%  §6.  Shadow Tables
%% =====================================================================

\section{Shadow Tables}\label{sec:shadow-tables}

All entries in the tables below are verified against the compute layer.
The modular characteristic $\kappa(\cA)$ is the arity-$2$ projection
of the universal MC element $\Theta_\cA$, and the genus-$g$ free energy
is $F_g(\cA) = \kappa(\cA) \cdot \lambda_g^{\mathrm{FP}}$, where
$\lambda_g^{\mathrm{FP}} = \frac{2^{2g-1}-1}{2^{2g-1}}
\cdot \frac{|B_{2g}|}{(2g)!}$
are the Faber--Pandharipande intersection numbers.

\subsection{Master shadow table}\label{ssec:master-shadow-table}

Table~\ref{tab:master-shadow} records, for each standard family,
the modular characteristic $\kappa$, the shadow class, the shadow
depth $r_{\max}$, the cubic coefficient $\alpha = S_3$, the quartic
coefficient $S_4$, and the critical discriminant $\Delta = 8\kappa S_4$.

\begin{table}[ht]
\centering
\caption{Master shadow table: standard landscape}
\label{tab:master-shadow}
\renewcommand{\arraystretch}{1.35}
{\small
\begin{tabular}{lcccccc}
\toprule
\textbf{Algebra $\cA$}
  & $\boldsymbol{\kappa(\cA)}$
  & \textbf{Class}
  & $\boldsymbol{r_{\max}}$
  & $\boldsymbol{S_3}$
  & $\boldsymbol{S_4}$
  & $\boldsymbol{\Delta}$ \\
\midrule
\multicolumn{7}{l}{\textit{Class G (Gaussian)}} \\[2pt]
Heisenberg $\cH_k$
  & $k$
  & G & $2$ & $0$ & $0$ & $0$ \\
Lattice $V_\Lambda$
  & $\operatorname{rank}(\Lambda)$
  & G & $2$ & $0$ & $0$ & $0$ \\
Free fermion $\psi$
  & $1/4$
  & G & $2$ & $0$ & $0$ & $0$ \\
$bc$ ghosts (weight $\lambda$)
  & $-(6\lambda^2{-}6\lambda{+}1)$
  & G & $2$ & $0$ & $0$ & $0$ \\[3pt]
\midrule
\multicolumn{7}{l}{\textit{Class L (Lie/tree)}} \\[2pt]
$\widehat{\mathfrak{sl}}_2$ at level $k$
  & $3(k{+}2)/4$
  & L & $3$
  & $\neq 0$ & $0$ & $0$ \\
$\widehat{\mathfrak{sl}}_N$ at level $k$
  & $(N^2{-}1)(k{+}N)/(2N)$
  & L & $3$
  & $\neq 0$ & $0$ & $0$ \\
$\widehat{E}_6$ at level $k$
  & $13(k{+}12)/4$
  & L & $3$
  & $\neq 0$ & $0$ & $0$ \\
$\widehat{E}_7$ at level $k$
  & $133(k{+}18)/36$
  & L & $3$
  & $\neq 0$ & $0$ & $0$ \\
$\widehat{E}_8$ at level $k$
  & $62(k{+}30)/15$
  & L & $3$
  & $\neq 0$ & $0$ & $0$ \\[3pt]
\midrule
\multicolumn{7}{l}{\textit{Class C (Contact)}} \\[2pt]
$\beta\gamma$ (weight $\lambda$)
  & $6\lambda^2{-}6\lambda{+}1$
  & C & $4$
  & $0$ & $0^\dagger$ & $0^\dagger$ \\[3pt]
\midrule
\multicolumn{7}{l}{\textit{Class M (Mixed)}} \\[2pt]
$\mathrm{Vir}_c$
  & $c/2$
  & M & $\infty$
  & $2$
  & $\dfrac{10}{c(5c{+}22)}$
  & $\dfrac{40}{5c{+}22}$ \\[8pt]
$\cW_3$ at $c$
  & $5c/6$
  & M & $\infty$
  & $\neq 0$ & $\neq 0$ & $\neq 0$ \\
$\cW_4$ at $c$
  & $13c/12$
  & M & $\infty$
  & $\neq 0$ & $\neq 0$ & $\neq 0$ \\
$\cW_5$ at $c$
  & $77c/60$
  & M & $\infty$
  & $\neq 0$ & $\neq 0$ & $\neq 0$ \\
$\cW_N$ at $c$ (general)
  & $(H_N{-}1)\,c$
  & M & $\infty$
  & $\neq 0$ & $\neq 0$ & $\neq 0$ \\
\bottomrule
\end{tabular}
}
\end{table}

\noindent
${}^\dagger$On the primary line, $S_4 = 0$ and $\Delta = 0$
for $\beta\gamma$; the quartic contact invariant
$Q^{\mathrm{contact}} \neq 0$ lives on a charged stratum
(stratum separation, Definition~\ref{def:class-C}).

\noindent
In the table, $H_N = 1 + 1/2 + \dotsb + 1/N$ denotes the
$N$-th harmonic number.  The Virasoro quartic contact
$S_4 = Q^{\mathrm{contact}}_{\mathrm{Vir}} = 10/[c(5c+22)]$
is the unique quartic shadow on the $T$-line, proved by direct
computation from the Virasoro OPE.  The T-line shadow data ($S_3 = 2$,
$S_4 = 10/[c(5c+22)]$) is \emph{universal}: it coincides for $\cW_N$
at all $N \geq 2$, since every $\cW_N$ contains the Virasoro
subalgebra along the stress-tensor primary line.  The additional
primary lines (spin-$3$, spin-$4$, etc.) carry independent shadow
data with separate quartic coefficients.


\subsection{Exceptional and non-simply-laced types}\label{ssec:exceptional-table}

Table~\ref{tab:exceptional-shadow} records the affine Kac--Moody
data for exceptional and non-simply-laced types.  All are class~L.

\begin{table}[ht]
\centering
\caption{Shadow data for exceptional and non-simply-laced types}
\label{tab:exceptional-shadow}
\renewcommand{\arraystretch}{1.35}
{\small
\begin{tabular}{lcccccc}
\toprule
$\fg$
  & $\dim(\fg)$
  & $h^\vee$
  & $\boldsymbol{\kappa\bigl(V_k(\fg)\bigr)}$
  & $c(V_k) + c(V_{k'})$
  & \textbf{Class}
  & $r_{\max}$ \\
\midrule
$B_2 \cong \mathfrak{so}_5$
  & $10$ & $3$ & $5(k{+}3)/3$ & $20$ & L & $3$ \\
$C_2 \cong \mathfrak{sp}_4$
  & $10$ & $3$ & $5(k{+}3)/3$ & $20$ & L & $3$ \\
$G_2$
  & $14$ & $4$ & $7(k{+}4)/4$ & $28$ & L & $3$ \\
$F_4$
  & $52$ & $9$ & $26(k{+}9)/9$ & $104$ & L & $3$ \\
$E_6$
  & $78$ & $12$ & $13(k{+}12)/4$ & $156$ & L & $3$ \\
$E_7$
  & $133$ & $18$ & $133(k{+}18)/36$ & $266$ & L & $3$ \\
$E_8$
  & $248$ & $30$ & $62(k{+}30)/15$ & $496$ & L & $3$ \\
\bottomrule
\end{tabular}
}
\end{table}

\noindent
The kappa formula is universal:
$\kappa(V_k(\fg)) = \dim(\fg)\,(k + h^\vee)/(2h^\vee)$.
The central charge sum $c + c' = 2\dim(\fg)$ under the
Feigin--Frenkel involution $k \mapsto -k - 2h^\vee$, and
the complementarity $\kappa + \kappa' = 0$ holds for all
affine Kac--Moody algebras.  For non-simply-laced types,
the dual Coxeter number $h^\vee$ differs from the Coxeter
number $h$ ($h^\vee = 3$ for $B_2$ and $C_2$, while $h = 4$).
The shadow class is~L regardless: the Jacobi identity kills the
quartic on every primary line.

\begin{remark}[$B_2$/$C_2$ isomorphism]
The Lie algebras $B_2 = \mathfrak{so}_5$ and
$C_2 = \mathfrak{sp}_4$ are isomorphic, sharing $\dim = 10$,
$h^\vee = 3$, and $h = 4$.  Their Cartan matrices are transposes
of each other, but all shadow invariants coincide.
\end{remark}


\subsection{Complementarity table}\label{ssec:complementarity-table}

\begin{table}[ht]
\centering
\caption{Complementarity data: $\kappa(\cA) + \kappa(\cA^!)$ and
$c(\cA) + c(\cA^!)$ under Koszul duality}
\label{tab:complementarity}
\renewcommand{\arraystretch}{1.35}
{\small
\begin{tabular}{llccc}
\toprule
\textbf{Family $\cA$}
  & \textbf{Koszul dual $\cA^!$}
  & $\boldsymbol{c + c'}$
  & $\boldsymbol{\kappa + \kappa'}$
  & \textbf{Self-dual point} \\
\midrule
$\cH_k$
  & $\mathrm{Sym}^{\mathrm{ch}}(V^*)$
  & --- & $0$ & none \\
$V_\Lambda$ (lattice)
  & $V_{\Lambda^!}$
  & --- & $0$ & $\Lambda$ self-dual \\
Free fermion
  & $\mathrm{Sym}^{\mathrm{ch}}(\gamma)$
  & $0$ & $0$ & --- \\
$\beta\gamma_\lambda$
  & $bc_\lambda$
  & $0$ & $0$ & --- \\[3pt]
\midrule
$V_k(\mathfrak{sl}_2)$
  & $V_{-k-4}(\mathfrak{sl}_2)$
  & $6$ & $0$ & $k = -2$ (critical) \\
$V_k(\mathfrak{sl}_N)$
  & $V_{-k-2N}(\mathfrak{sl}_N)$
  & $2(N^2{-}1)$ & $0$ & $k = -N$ (critical) \\
$V_k(\fg)$ (general)
  & $V_{-k-2h^\vee}(\fg)$
  & $2\dim(\fg)$ & $0$ & $k = -h^\vee$ (critical) \\[3pt]
\midrule
$\mathrm{Vir}_c$
  & $\mathrm{Vir}_{26-c}$
  & $26$ & $13$ & $c = 13$ \\
$\cW_3$ at $c$
  & $\cW_3$ at $100{-}c$
  & $100$ & $250/3$ & $c = 50$ \\
$\cW_4$ at $c$
  & $\cW_4$ at $246{-}c$
  & $246$ & $533/2$ & $c = 123$ \\
$\cW_5$ at $c$
  & $\cW_5$ at $488{-}c$
  & $488$ & $9394/15$ & $c = 244$ \\
$\cW_N$ at $c$ (general)
  & $\cW_N$ at $K_N{-}c$
  & $K_N$ & $(H_N{-}1)\,K_N$ & $c = K_N/2$ \\
\bottomrule
\end{tabular}
}
\end{table}

\noindent
Here $K_N = 2(N{-}1) + 4(N^2{-}1)N$ is the Koszul conductor for
$\cW_N = \cW(\mathfrak{sl}_N)$ (the sum $c + c'$ under the
Feigin--Frenkel involution $k \mapsto -k - 2N$ on the underlying
$\mathfrak{sl}_N$).  The key structural distinction: for all
affine Kac--Moody algebras, $\kappa + \kappa' = 0$ (anti-symmetry
under the Feigin--Frenkel involution); for $\cW$-algebras,
$\kappa + \kappa' = (H_N - 1) \cdot K_N \neq 0$.
The Virasoro ($N = 2$) gives $\kappa + \kappa' = \tfrac{1}{2} \cdot 26
= 13$, which vanishes only at the self-dual point $c = 13$
where $\kappa(\mathrm{Vir}_{13}) = \kappa(\mathrm{Vir}_{13}) = 13/2$.


\subsection{Genus tables}\label{ssec:genus-tables}

The Faber--Pandharipande numbers and the genus expansion
$F_g(\cA) = \kappa(\cA) \cdot \lambda_g^{\mathrm{FP}}$.

\begin{table}[ht]
\centering
\caption{Faber--Pandharipande intersection numbers $\lambda_g^{\mathrm{FP}}$}
\label{tab:faber-pandharipande}
\renewcommand{\arraystretch}{1.35}
{\small
\begin{tabular}{ccc}
\toprule
$g$ & $\lambda_g^{\mathrm{FP}}$ & Decimal \\
\midrule
$1$ & $1/24$ & $0.041\overline{6}$ \\
$2$ & $7/5760$ & $1.215 \times 10^{-3}$ \\
$3$ & $31/967680$ & $3.204 \times 10^{-5}$ \\
$4$ & $127/154828800$ & $8.201 \times 10^{-7}$ \\
$5$ & $73/3503554560$ & $2.084 \times 10^{-8}$ \\
\bottomrule
\end{tabular}
}
\end{table}

\noindent
The formula is
$\lambda_g^{\mathrm{FP}} = \frac{2^{2g-1} - 1}{2^{2g-1}}
\cdot \frac{|B_{2g}|}{(2g)!}$,
where $B_{2g}$ is the $2g$-th Bernoulli number.  The generating
function is $\sum_{g \geq 1} \lambda_g^{\mathrm{FP}}\, \hbar^{2g}
= \hat{A}(i\hbar) - 1$, where $\hat{A}(x) = (x/2)/\sin(x/2)$ is
the $\hat{A}$-genus.  The numbers $\lambda_g^{\mathrm{FP}}$ are
strictly positive, strictly decreasing, and decay like
$(2\pi)^{-2g}$.

\begin{table}[ht]
\centering
\caption{Genus expansion $F_g(\cA) = \kappa \cdot \lambda_g^{\mathrm{FP}}$
for selected families}
\label{tab:genus-expansion}
\renewcommand{\arraystretch}{1.35}
{\small
\begin{tabular}{lccccc}
\toprule
\textbf{Algebra $\cA$}
  & $\boldsymbol{\kappa}$
  & $\boldsymbol{F_1}$
  & $\boldsymbol{F_2}$
  & $\boldsymbol{F_3}$ \\
\midrule
$\cH_1$ (Heisenberg, $k{=}1$)
  & $1$ & $1/24$ & $7/5760$ & $31/967680$ \\
$V_1(\mathfrak{sl}_2)$
  & $9/4$ & $3/32$ & $7/2560$ & $31/430080$ \\
$V_1(\mathfrak{sl}_3)$
  & $16/3$ & $2/9$ & $7/1080$ & $31/181440$ \\
$V_\Lambda$ (Leech, $\operatorname{rank}{=}24$)
  & $24$ & $1$ & $7/240$ & $31/40320$ \\
$\beta\gamma$ (standard, $\lambda{=}1$)
  & $1$ & $1/24$ & $7/5760$ & $31/967680$ \\
$\mathrm{Vir}_1$
  & $1/2$ & $1/48$ & $7/11520$ & $31/1935360$ \\
$\mathrm{Vir}_{13}$
  & $13/2$ & $13/48$ & $91/11520$ & $403/1935360$ \\
$\mathrm{Vir}_{26}$
  & $13$ & $13/24$ & $91/5760$ & $403/967680$ \\
$\cW_3$ at $c{=}6$
  & $5$ & $5/24$ & $7/1152$ & $31/193536$ \\
\bottomrule
\end{tabular}
}
\end{table}

\noindent
The entries for $V_1(\mathfrak{sl}_3)$ use
$\kappa = (8 \cdot 4)/(2 \cdot 3) = 16/3$.  The Leech lattice
($\operatorname{rank} = 24$) gives $F_1 = 1$, reflecting the
Jacobi $\Delta$-function origin of the Leech theta series.

\subsection{Shadow growth rate table}\label{ssec:growth-rate-table}

\begin{table}[ht]
\centering
\caption{Shadow growth rate $\rho(\cA)$ for class~M algebras
(Virasoro at selected central charges)}
\label{tab:growth-rate}
\renewcommand{\arraystretch}{1.35}
{\small
\begin{tabular}{ccccl}
\toprule
$c$
  & $\kappa$
  & $Q^{\mathrm{contact}}$
  & $\rho$
  & \textbf{Regime} \\
\midrule
$1/2$ (Ising)
  & $1/4$ & $40/49$ & $\approx 12.53$ & divergent \\
$1$ (free boson)
  & $1/2$ & $10/27$ & $\approx 6.24$ & divergent \\
$4$
  & $2$ & $5/84$ & $\approx 1.54$ & divergent \\
$c^* \approx 6.125$
  & $\approx 3.06$ & $\approx 0.031$ & $1$ & critical \\
$13$ (self-dual)
  & $13/2$ & $10/1131$ & $\approx 0.467$ & convergent \\
$25$
  & $25/2$ & $2/735$ & $\approx 0.242$ & convergent \\
$26$ (bosonic string)
  & $13$ & $5/1976$ & $\approx 0.232$ & convergent \\
\bottomrule
\end{tabular}
}
\end{table}

\noindent
The shadow growth rate $\rho = \sqrt{(9\alpha^2 + 2\Delta)/(4\kappa^2)}$
determines the convergence of the shadow tower.  For Virasoro,
$\rho^2 = (180c + 872)/[c^2(5c+22)]$.  The critical central charge
$c^* \approx 6.125$ (the positive real root of~\eqref{eq:critical-cubic})
separates the convergent regime ($c > c^*$, tower summable) from the
divergent regime ($c < c^*$, tower grows exponentially).  Class~G and
class~L algebras have $\rho = 0$ (the tower terminates, so there is no
exponential growth).

\subsection{Shadow field table}\label{ssec:shadow-field-table}

For class~$\mathbf{M}$ algebras at rational~$c$, the shadow metric
defines an imaginary quadratic field $K_L = \mathbb{Q}(\sqrt{D_0})$
through which the Epstein zeta factors.  The discriminant
$\operatorname{disc}(Q_L) = -32\kappa^2\Delta$ is negative for
positive~$c$.

\begin{table}[ht]
\centering
\caption{Shadow fields for the Virasoro algebra at rational~$c$.
The squarefree kernel $D_0$ is computed from
$\operatorname{disc}(Q_L) = -320c^2/(5c+22)$;
$K_L = \mathbb{Q}(\sqrt{D_0})$.}
\label{tab:shadow-fields-vir}
\renewcommand{\arraystretch}{1.35}
{\small
\begin{tabular}{lcccccl}
\toprule
\textbf{Model} & $c$ & $\kappa$ & $\Delta$
  & $D_0$
  & $h(D_0)$
  & \textbf{Shadow field} \\
\midrule
Ising $\mathcal{M}(3,4)$ & $1/2$ & $1/4$ & $80/49$
  & $-10$ & $2$
  & $\mathbb{Q}(\sqrt{-10})$ \\
Tricritical Ising & $7/10$ & $7/20$ & $200/309$
  & $-510$ & ---
  & $\mathbb{Q}(\sqrt{-510})$ \\
Free boson & $1$ & $1/2$ & $40/27$
  & $-15$ & $2$
  & $\mathbb{Q}(\sqrt{-15})$ \\
$c = 4$ & $4$ & $2$ & $20/21$
  & $-210$ & ---
  & $\mathbb{Q}(\sqrt{-210})$ \\
$c^* \approx 6.12$ & $c^*$ & $\approx 3.06$ & $\approx 0.99$
  & ---  & ---
  & (critical) \\
Self-dual & $13$ & $13/2$ & $40/87$
  & $-435$ & ---
  & $\mathbb{Q}(\sqrt{-435})$ \\
Critical string & $26$ & $13$ & $40/152$
  & $-190$ & ---
  & $\mathbb{Q}(\sqrt{-190})$ \\
\bottomrule
\end{tabular}
}
\end{table}

\noindent
Here $D_0$ is the squarefree kernel of the discriminant,
$h(D_0)$ is the class number of~$\mathbb{Q}(\sqrt{D_0})$,
and the discriminant formula is
$\operatorname{disc}(Q_L) = -320\,c^2/(5c+22)$, so
$D_0 = \operatorname{sqfree}(-5\,q\,(5p + 22q))$ for
$c = p/q$ in lowest terms.  The Ising model ($D_0 = -10$)
and the free boson ($D_0 = -15$) both have class number~$2$;
by Davenport--Heilbronn, the corresponding Epstein zeta
functions may have zeros off the critical line.


\subsection{Virasoro shadow tower table}\label{ssec:virasoro-tower-table}

\begin{table}[ht]
\centering
\caption{Virasoro shadow coefficients $S_r$ for $r = 2, \dotsc, 10$}
\label{tab:virasoro-tower}
\renewcommand{\arraystretch}{1.35}
{\small
\begin{tabular}{cl}
\toprule
$r$ & $S_r(\mathrm{Vir}_c)$ \\
\midrule
$2$ & $c/2$ \\
$3$ & $2$ \\
$4$ & $10/[c(5c+22)]$ \\
$5$ & $-48/[c^2(5c+22)]$ \\
$6$ & $80(45c+193)/[3c^3(5c+22)^2]$ \\
$7$ & $-2880(15c+61)/[7c^4(5c+22)^2]$ \\
$8$ & $80(2025c^2+16470c+33314)/[c^5(5c+22)^3]$ \\
$9$ & $-1280(2025c^2+15570c+29554)/[3c^6(5c+22)^3]$ \\
$10$ & $256(91125c^3+1050975c^2+3989790c+4969967)/[c^7(5c+22)^4]$ \\
\bottomrule
\end{tabular}
}
\end{table}

\noindent
Each $S_r$ is a rational function of $c$ with poles at $c = 0$ and
$c = -22/5$; these are the two values where the quadratic Casimir
$c(5c+22)$ vanishes.  The sign pattern
$S_r > 0$ for even $r$ and $S_r < 0$ for odd $r \geq 5$
(at $c > 0$) reflects the oscillatory decay $\cos(r\theta + \phi)$
from the complex branch points of the shadow generating function.
All entries have been verified by two independent methods: closed-form
derivation and the convolution recursion from $H^2 = t^4 Q_L$.
